
\subsection*{Abstract}
The "Black-Box" nature of modern Graph Neural Networks (GNNs) and Transformer-based reasoning systems limits their utility in domains requiring high auditability. We present the \textbf{Glass-Box Reasoning Studio}, an interactive visualization framework designed for the forensic audit of multi-hop Knowledge Graph inference. The Studio reifies the "Reasoning Beam" as a dynamic topological trace, where edges are scaled by their \textbf{Community-Structured Attention (CSA)} weights and nodes are color-coded by their \textbf{DSCF/TSC} community partitions. We introduce a "Forensic Score Breakdown" interface that exposes the latent mathematical signals (semantic similarity, community guidance, and structural centrality) driving each traversal hop, building on foundational Explainable AI (XAI) principles \cite{samek2017explainable, ribeiro2016lime, lundberg2017shap}. Furthermore, we describe a real-time "Live Feed" visualization for streaming graphs that animates \textbf{STDP spike events} and the materialization of speculative causal links. The v2.71.0 release adds adaptive node clustering to support visual scaling for graphs exceeding $10^5$ nodes. In v2.71.0 (Phase 54), a major architectural refactor extracts all Studio business logic into \texttt{core/studio\_engine.py} (StudioEngine class), enabling 38 new unit tests that run without a live Gradio server. The 10-parameter CSA weight profiler is corrected to expose all parameters including $\mu$ (synthesis-density penalty), and a new dark-mode monitoring dashboard with live log streaming is introduced. Our results show that this interactive "Glass-Box" approach significantly reduces the time required for human experts to verify complex AI-generated hypotheses.

\subsection*{1. Introduction}
Explainability in AI (XAI) has traditionally focused on post-hoc interpretations of neural weights (e.g., saliency maps). In graph-based reasoning, however, the explanation is the path itself. The Glass-Box Reasoning Studio provides the first integrated environment for visualizing graph attention as a physical, navigatable flow.

\subsection*{2. Forensic Visualization Methodology}

\subsubsection*{2.1 The Reasoning Trace}
The Studio implements a path-centric rendering algorithm that isolates the sub-graph involved in a specific query. The attention weight $a(u,v,k)$ is mapped to edge thickness and opacity, allowing the user to visually perceive the "narrowing of the beam" as the AI focuses on likely answers.

\subsubsection*{2.2 Modal Animations}
For temporal and streaming data, the Studio utilizes high-frequency state updates:
\begin{itemize}
  \item \textbf{Potentiation}: Edges being strengthened by LTP [Buchorn, 2026] increase in saturation.
  \item \textbf{Drift}: Community boundaries shift smoothly using force-directed layouts to reflect modularity updates [Buchorn, 2026].
\end{itemize}


\subsection*{3. Interactive Debugging (v2.71.0)}
The Studio provides a "Dialectical reasoning" mode where users can manually adjust CSA parameters ($\alpha, \beta, \gamma$) via sliders and observe the immediate physical shift in the reasoning beam, providing a "Human-in-the-Loop" (HITL) interface for hyperparameter tuning. In v2.71.0, this includes real-time feedback submission to the \textbf{MetaParameterLearner}.

\subsection*{4. Recent Advances (v2.51.1 -> v2.71.0)}

\subsubsection*{4.1 StudioEngine Architectural Refactor (Phase 54)}
The most significant change in v2.71.0 is a complete architectural separation of Studio business logic from the Gradio server layer. Previously, all reasoning coordination, graph management, and query dispatch were embedded directly in \texttt{ui/studio.py}. Phase 54 extracts these into a new \texttt{core/studio\_engine.py} module exposing the \texttt{StudioEngine} class.

Benefits of this separation:
\begin{itemize}
  \item \textbf{Independent testability}: 38 new unit tests exercise StudioEngine directly without requiring a running Gradio server, reducing test fragility and CI/CD runtime.
  \item \textbf{Reusability}: StudioEngine can be instantiated by the REST API server (\texttt{api/server.py}) and the CLI without the UI layer.
  \item \textbf{Separation of concerns}: \texttt{ui/studio.py} is reduced to a thin Gradio binding layer; all algorithmic logic lives in \texttt{core/}.
\end{itemize}


\subsubsection*{4.2 10-Parameter CSA Weight Profiler (Bug Fix)}
The Studio's CSA weight profiler previously exposed only 9 of the 10 CSA parameters, omitting $\mu$ (synthesis-density penalty). This meant that Studio users tuning parameters interactively could not adjust the penalty applied to paths over-relying on Synaptic Bridge-synthesized edges. In v2.71.0, the profiler exposes all 10 parameters:

\begin{table}[h]
\small\centering
\begin{tabular}{lll}
\toprule
Parameter & Symbol & Description \\
\midrule
alpha & $\alpha$ & Semantic similarity (cosine) \\
beta & $\beta$ & Community score \\
gamma & $\gamma$ & Edge-type weight \\
delta & $\delta$ & Normalized distance penalty \\
epsilon & $\varepsilon$ & Hop decay \\
zeta & $\zeta$ & PageRank prior \\
eta & $\eta$ & Temporal decay \\
iota & $\iota$ & Node recency \\
\textbf{mu} & \textbf{$\mu$} & \textbf{Synthesis-density penalty (was missing)} \\
theta & $\theta$ & Grounding confidence \\
\bottomrule
\end{tabular}
\end{table}


\subsubsection*{4.3 Hot CSV Reload via /build Endpoint (Phase 54)}
A new \texttt{/build} REST endpoint accepts a CSV file upload and triggers a live graph rebuild without server restart. Studio users can iterate on graph construction --- adding new entities, edges, or relation types --- and immediately see the updated reasoning behavior in the visualization layer. This enables rapid prototyping workflows where graph and query patterns are co-developed.

\subsubsection*{4.4 Dark-Mode Monitoring Dashboard (dashboard.html)}
A new \texttt{ui/dashboard.html} provides a production monitoring interface built on GridStack (resizable widget layout), Chart.js (time-series metrics), and vis-network (live graph visualization). Key panels:
\begin{itemize}
  \item \textbf{Live query throughput} (queries/sec, rolling 60s window)
  \item \textbf{CSA parameter drift} (time-series of all 10 parameters as MetaParameterLearner updates them)
  \item \textbf{Community partition map} (vis-network rendering, auto-refreshed on rebalance)
  \item \textbf{Log stream} (real-time display from \texttt{/logs} ring buffer)
\end{itemize}


\subsubsection*{4.5 /logs Endpoint with Ring Buffer}
The Studio and API server now expose a \texttt{/logs} GET endpoint backed by a \texttt{RingBufferHandler} that captures all \texttt{cerebrum.*} log events at DEBUG level. The ring buffer holds the last N log entries (configurable, default 1,000) and returns them as structured JSON. A DELETE on \texttt{/logs} clears the buffer. This enables the monitoring dashboard to display live operational state without requiring external logging infrastructure.

\subsection*{5. Conclusion}
The Glass-Box Reasoning Studio transforms graph attention from an abstract mathematical construct into a tangible, auditable artifact. In v2.71.0, the Phase 54 architectural refactor (StudioEngine extraction, 38 new unit tests), the corrected 10-parameter weight profiler, the \texttt{/build} hot-reload endpoint, and the dark-mode monitoring dashboard collectively advance the Studio from an interactive demo into a production-grade reasoning observatory. By bridging the gap between latent semantic operations and human-readable topologies, it enables the deployment of autonomous reasoning systems in high-stakes environments.